\subsection{Weakly complete spaces}

PROPOSITION 9. --- Let F, G be two vector spaces in separating duality. If \(\hat{\mathbf{F}}\) is the completion of the space F for the topology \(\sigma(\mathrm{F},\mathrm{G})\) and if we consider the canonical injection \(j:\mathrm{F}\to\mathrm{G}^{*}\) , where \(\mathrm{G}^{*}\) has the topology \(\sigma(\mathrm{G}^{*},\mathrm{G})\) , then the continuous extension \(\hat{j}:\hat{\mathrm{F}}\to\mathrm{G}^{*}\) of \(j\) is an isomorphism of topological vector spaces.

For, we see that \(G^{*}\) , endowed with \(\sigma(G^{*}, G)\) , is Hausdorff and complete (II, p.~51, cor. 2); if we identify F by j with a vector subspace of \(G^{*}\) then the topology induced on F by \(\sigma(G^{*}, G)\) is \(\sigma(F, G)\) , and F is dense in \(G^{*}\) in the topology \(\sigma(G^{*}, G)\) (II, p.~43, cor. 4); from which the proposition follows.

Vector spaces that are complete for a weak topology are therefore the duals G* of arbitrary vector spaces G endowed with \(\sigma(\mathbf{G}^{*}, \mathbf{G})\) ; after II, p.~51, cor. 2, they are (topologically) isomorphic to products \(R^{I}\) of real lines. To simplify the language, we shall call them products of lines (for an intrinsic characterisation of these spaces see II, p.~85, exerc. 13 and II, p.~81, exerc. 1).

We note that on \(G^{*}\) , the \(\sigma(G^{*}, G)\) topology is minimal among the weak topologies that are Hausdorff; for, a weak topology that is coarser than \(\sigma(G^{*}, G)\) is necessarily of the form \(\sigma(G^{*}, H)\) where \(H \subset G\) (II, p.~43, cor. 3); but if \(H \neq G\) , then there exists a linear form \(x^{*} \in \mathbf{G}^{*}\) that is non null and is orthogonal to \(\mathbf{H}\) (A, II, § 7.3, prop. 8), therefore \(\sigma(\mathbf{G}^{*}, \mathbf{H})\) is not Hausdorff.

We deduce from this remark that, if F, G are two vector spaces, then a linear bijection \(u:G^{*}\to F^{*}\) , that is continuous for the topologies \(\sigma(G^{*},G)\) and \(\sigma(F^{*},F)\) , is necessarily bicontinuous.

PROPOSITION 10. --- Let G be a real vector space and F = G* its dual with the topology \(\sigma(\mathbf{G}^*, \mathbf{G})\) .

\begin{enumerate}
\def\labelenumi{(\roman{enumi})}
\item
  The mapping \(\mathbf{V} \mapsto \mathbf{V}^{\circ}\) is a bijection of the set of vector subspaces of \(\mathbf{G}\) on the set of closed vector subspaces of \(\mathbf{F}\) .
\item
  Every closed vector subspace of \(\mathbf{F}\) is a product of lines and has a topological complement.
\end{enumerate}

By the bipolar theorem (II, p.~45, cor. 2) \(\mathrm{V} \mapsto \mathrm{V}^{\circ}\) is a bijection of the set of vector subspaces V of G, closed for \(\sigma(\mathbf{G}, \mathbf{G}^{*})\) on the set of closed vector subspaces of F. But, by definition, every linear form on G is continuous for \(\sigma(\mathbf{G}, \mathbf{G}^{*})\) , therefore every vector subspace in G is closed, being defined by a system of equations \(\langle y, y_{\lambda}^{*} \rangle = 0\) (where \(y_{\lambda}^{*} \in \mathbf{G}^{*}\) ); this proves (i).

Now let W be a closed subspace of F; we have then \(W = V^{\circ}\) with \(V = W^{\circ}\) in G. Let \(V'\) be a complement of V in G. We know that \(F = G^{*}\) can be canonically identified with \(V^{*} \oplus V'^{*}\) , and \(V'^{*}\) identified with \(V^{\circ} = W\) (A, II, § 2.6, cor. to prop. 10); further (II, p.~50, prop. 8) the topology \(\sigma(G^{*}, G)\) can be identified with the product of the topologies \(\sigma(V^{*}, V)\) and \(\sigma(V'^{*}, V')\) ; this proves assertion (ii).

Though, for the topology \(\sigma(\mathbf{G},\mathbf{G}^{*})\) , every vector subspace of \(\mathbf{G}\) is closed, we note that if \(\mathbf{G}\) is of infinite dimension then the topology \(\sigma(\mathbf{G},\mathbf{G}^{*})\) is not the finest locally convex topology on \(\mathbf{G}\) , every neighbourhood of 0 for \(\sigma(\mathbf{G},\mathbf{G}^{*})\) containing a vector subspace of infinite dimension: it is however the finest of the weak topologies on \(\mathbf{G}\) (II, p.~43, cor. 3).

